%  03_dataset.tex
\section{Dataset}
\label{sec:dataset}

\subsection{Species and capture conditions}

We work with two fruit-bat species from active research colonies:

\begin{itemize}[itemsep=2pt]
  \item \textbf{\rousettus (Egyptian fruit bat)}: 12 individuals, 1\,059 images.
        Images were extracted from close-range video recordings; each frame was
        independently processed by the preprocessing pipeline.  All images were captured
        with randomised background compositing.
  \item \textbf{\mauritius (Mauritian flying fox, \emph{Pteropus niger})}: three
        background conditions are available for this species, each captured/curated
        independently and therefore differing in identity and image counts:
        \emph{random} (composited random textures; 11 individuals, 1\,119 images),
        \emph{green screen} (chroma-key replacement; 19 individuals, 1\,598 images),
        and \emph{original} (natural enclosure background; 17 individuals, 1\,038
        images).  Together they provide a controlled study of background-invariance
        (see Table~\ref{tab:dataset} for per-split counts).
\end{itemize}

\subsection{Preprocessing pipeline}

Raw video frames undergo a five-stage preprocessing pipeline before storage.
First, a YOLO-v8 pose estimator (custom weights \texttt{face\_pose.pt}) detects facial
keypoints; the face is then rotated so that the eyes are horizontally aligned, providing
pose-normalised crops.  A second YOLO-v8 model (\texttt{face\_seg.pt}) segments the bat
face region; the bounding box of the predicted mask drives a square crop with a
configurable outer margin (default $10\%$).  Optional background replacement composites
one of several procedurally generated backgrounds (blur, gradient, random noise, or
downloaded textures).  Finally, images are resized to the model's input resolution and
normalised to $[0, 1]$.

Image \emph{quality} is quantified per frame as the variance of the Laplacian
$\sigma^2_\Delta = \mathrm{Var}(\Delta I)$, a standard focus measure; this score is stored
in the manifest and can be used to filter low-quality frames during dataset construction.

\subsection{Manifest structure}

All dataset metadata is stored in a \emph{manifest}: a CSV file with the following
columns: \texttt{path}, \texttt{identity}, \texttt{species}, \texttt{background},
\texttt{source}, \texttt{augmented}, \texttt{split}, \texttt{quality}.
Each manifest is accompanied by a SHA-256 hash sidecar; loading a manifest with a
mismatched hash raises an error, preventing silent data-version drift.

\subsection{Identity-disjoint splitting}

The central methodological contribution of this work is the enforcement of
\emph{identity-disjoint} splits.  Given a set of $N$ unique individual identities,
the splitter assigns each identity to exactly one of \{train, val, test\} before any
images are allocated.  This guarantees that the model never observes images of a test
identity during training or validation---the protocol required for honest open-set
evaluation.

Formally, let $\mathcal{I}$ be the set of identities.  The splitter draws
$(n_{\mathrm{val}}, n_{\mathrm{test}})$ identities as
$n_{\mathrm{val}} = \lfloor |\mathcal{I}| \cdot r_{\mathrm{val}} \rceil$,
$n_{\mathrm{test}} = \lfloor |\mathcal{I}| \cdot r_{\mathrm{test}} \rceil$,
with the remainder assigned to training.  The shuffle is seeded (seed 42) over an
alphabetically sorted identity list, making the split fully deterministic.

Figure~\ref{fig:split_comparison} contrasts the random-pair protocol with the
identity-disjoint protocol.

\begin{figure}[t]
\centering
\begin{subfigure}[t]{0.46\linewidth}
\centering
\begin{tikzpicture}[
  idnode/.style={draw, circle, minimum size=0.52cm, font=\small\bfseries, inner sep=0pt,
                 fill=gray!12},
  splitbox/.style={draw, rounded corners=3pt, minimum width=1.8cm, minimum height=0.5cm,
                   font=\scriptsize, align=center}
]
  % identities column
  \node[idnode] (a) at (0, 1.5) {A};
  \node[idnode] (b) at (0, 0.75) {B};
  \node[idnode] (c) at (0, 0)   {C};

  % split boxes
  \node[splitbox, fill=colSiamese!20, draw=colSiamese!60]
    (tr) at (2.6, 1.5) {Train};
  \node[splitbox, fill=colAdaFace!20, draw=colAdaFace!60]
    (va) at (2.6, 0.75) {Val};
  \node[splitbox, fill=colArcFace!20, draw=colArcFace!60]
    (te) at (2.6, 0) {Test};

  % every identity connects to every split (leaky)
  \foreach \src in {a, b, c} {
    \draw[->, gray!55, thin] (\src.east) to[out=5,in=175] (tr.west);
    \draw[->, gray!55, thin] (\src.east) to[out=0,in=180] (va.west);
    \draw[->, gray!55, thin] (\src.east) to[out=-5,in=185] (te.west);
  }

  \node[font=\tiny, red!70!black] at (1.3, -0.55) {\textbf{Leakage: all ids in all splits}};
\end{tikzpicture}
\caption{Random-pair split (leaky)}
\end{subfigure}
\hfill
\begin{subfigure}[t]{0.46\linewidth}
\centering
\begin{tikzpicture}[
  idnode/.style={draw, circle, minimum size=0.52cm, font=\small\bfseries, inner sep=0pt},
  splitbox/.style={draw, rounded corners=3pt, minimum width=1.8cm, minimum height=0.5cm,
                   font=\scriptsize, align=center}
]
  % identities column: 2 train, 1 val, 1 test
  \node[idnode, fill=colSiamese!30, draw=colSiamese] (a) at (0, 2.0) {A};
  \node[idnode, fill=colSiamese!30, draw=colSiamese] (b) at (0, 1.3) {B};
  \node[idnode, fill=colAdaFace!30, draw=colAdaFace](c) at (0, 0.6) {C};
  \node[idnode, fill=colArcFace!30, draw=colArcFace](d) at (0, -0.1){D};

  % split boxes
  \node[splitbox, fill=colSiamese!20, draw=colSiamese!60]
    (tr) at (2.6, 1.65) {Train};
  \node[splitbox, fill=colAdaFace!20, draw=colAdaFace!60]
    (va) at (2.6, 0.6) {Val};
  \node[splitbox, fill=colArcFace!20, draw=colArcFace!60]
    (te) at (2.6, -0.1) {Test};

  % each identity to exactly one split
  \draw[->, colSiamese, semithick] (a.east) to[out=0,in=180] (tr.west);
  \draw[->, colSiamese, semithick] (b.east) to[out=0,in=180] (tr.west);
  \draw[->, colAdaFace, semithick] (c.east) -- (va.west);
  \draw[->, colArcFace, semithick] (d.east) -- (te.west);

  \node[font=\tiny, colAdaFace!80!black] at (1.3, -0.65)
    {\textbf{No leakage: each id in one split only}};
\end{tikzpicture}
\caption{Identity-disjoint split (correct)}
\end{subfigure}
\caption{Splitting strategies.  (a)~The random-pair protocol assigns images without
regard to identity: every individual (A, B, C) contributes images to all three partitions.
(b)~The identity-disjoint protocol assigns each individual to exactly one partition, so
test identities (D) are never seen during training.}
\label{fig:split_comparison}
\end{figure}

\subsection{Dataset statistics}

Table~\ref{tab:dataset} summarises the final partition statistics.

\begin{table}[h]
\centering
\caption{Dataset split statistics under the identity-disjoint protocol (seed 42).
Whole identities are assigned to exactly one split; every condition reserves 3
identities for validation and 3 for test, with the remainder used for training.
BG~=~background condition.}
\label{tab:dataset}
\setlength{\tabcolsep}{5pt}
\begin{tabular}{llccccccc}
\toprule
Species & BG & Total & Identities & \multicolumn{2}{c}{Train} & \multicolumn{2}{c}{Val} & Test \\
\cmidrule(lr){5-6}\cmidrule(lr){7-8}
 & & imgs & total & ids & imgs & ids & imgs & imgs \\
\midrule
\rousettus       & random   & 1\,059 & 12 & 6  & 486   & 3 & 343 & 230 \\
\rousettus       & green    & 1\,093 & 12 & 6  & 490   & 3 & 359 & 244 \\
\mauritius       & random   & 1\,119 & 11 & 5  & 400   & 3 & 343 & 376 \\
\mauritius       & green    & 1\,598 & 19 & 13 & 1\,229& 3 & 230 & 139 \\
\mauritius       & original & 1\,038 & 17 & 11 & 675   & 3 & 175 & 188 \\
\bottomrule
\end{tabular}
\end{table}

The small identity counts---fewer than twelve training identities in all conditions---are
representative of typical ecological monitoring studies and define the challenge regime
analysed in Section~\ref{sec:experiments}.
